\section*{\texorpdfstring{\S\CurrentExerciseGroup}{Section \CurrentExerciseGroup}}
\phantomsection
\addcontentsline{toc}{section}{\texorpdfstring{Exercises for \S\CurrentExerciseGroup}{Exercises for Section \CurrentExerciseGroup}}

\begin{enumerate}
\def\labelenumi{\arabic{enumi})}
\item
  Let X be a locally compact space, \(\Phi\) a filter on \(\mathcal{M}(\mathrm{X};\mathbf{C})\) such that the support of \(\mu\) recedes indefinitely along \(\Phi\) (No.~2, preceding Prop. 7); show that for the quasi-strong topology ( \(\S1\) , Exer. 8), \(\mu\) tends to 0 with respect to \(\Phi\) .
\item
  Show that for the ultrastrong topology on \(\mathcal{M}^1 (\mathrm{X};\mathbf{C})\) (§1, Exer. 15), the set of measures with compact support is dense in \(\mathcal{M}^1 (\mathrm{X};\mathbf{C})\) .
\item
  Let \(\mathbf{X}\) be the interval {[}0,1{]} of \(\mathbf{R}\) . Show that in the Banach space \(\mathcal{M}(\mathbf{X};\mathbf{C})\) , the distance from the Lebesgue measure to every discrete measure is \(\geqslant 1\) .
\item
  Let A be an uncountable set. In the set \(\mathfrak{P}(A)\) , consider the relation \(\langle M \cap \complement N\) and \(N \cap \complement M\) are countable \(\rangle\) between M and N; show that this is an equivalence relation, weakly compatible with the inclusion relation \(M \subset M'\) (S, III, §1, Exer. 2). Let B be the quotient set of \(\mathfrak{P}(A)\) for this equivalence relation; show that on B, the relation deduced from the inclusion relation by passage to quotients is an order relation for which B is a Boolean algebra (loc. cit, Exer. 17). One knows (GT, II, §4, Exer. 12) that there exists an order structure isomorphism of B onto the Boolean algebra formed by the subsets that are both open and closed in a totally disconnected compact space X. Show that if U is a nonempty open set in X, then there exists an uncountable family of pairwise disjoint nonempty open subsets \(V_{\alpha}\) of U (make use of the fact that an uncountable set admits an uncountable partition consisting of uncountable sets). From this, deduce that every measure on X has nowhere dense support.
\item
  Let \(X\) be a compact space. An infinite sequence \((x_{n})_{n\in \mathbf{N}}\) of points of \(X\) is said
\end{enumerate}

to have a limit distribution if the sequence of measures with finite support \(\left(\sum_{i=0}^{n-1} \varepsilon_{x_i}\right)/n\) tends vaguely to a limit \(\mu\) in \(\mathcal{M}_{+}(\mathrm{X})\) ; one also says that the sequence \((x_{n})\) is equidistributed for the measure \(\mu\) . If \(\mathrm{T}\) is a total set in the Banach space \(\mathcal{C}(\mathrm{X};\mathbf{C})\) , show that for a sequence \((x_{n})\) to have a limit distribution, it is necessary and sufficient that the sequence \(\left(\left(\sum_{i=0}^{n-1} f(x_i)\right)/n\right)\) tend to a limit in \(\mathbf{C}\) for every function \(f \in \mathrm{T}\) . In particular, show that if \(X\) is a metrizable compact space, then for every sequence \((x_{n})\) of points of \(X\) there exists a subsequence \((x_{n_k})\) that has a limit distribution.

If X is the interval {[}0,1{]} of R and if \(\mu\) is Lebesgue measure on X, show that for a sequence \((x_{n})\) of points of X to be equidistributed for \(\mu\) , it is necessary and sufficient that it be equidistributed mod 1 in the sense of GT, VII, §1, Exer. 14. For this, it is necessary and sufficient that for every integer \(m \in Z\) , the sequence

\(\left(\left(\sum_{k=0}^{n-1} e^{2i\pi mx_k}\right)/n\right)\) tend to 0 for \(m \neq 0\) . From this deduce anew, in particular, that the sequence \(x_{n} = n\theta - [n\theta]\) ( \(n \in \mathbf{N}\) ) is equidistributed mod 1 for every irrational number \(\theta\) (GT, VII, §1, Exer. 14). What is its limit distribution when \(\theta\) is rational? Cf. also Ch. IV, §5, No.~12.

\setcounter{exercisegroup}{3}
\edef\CurrentExerciseGroup{\arabic{exercisegroup}}
